\documentclass[journal]{IEEEtran}

% Math packages
\usepackage{amsmath,amsfonts,amssymb,amsthm,bm,mathrsfs}

% Algorithm packages
\usepackage{algorithmic,algorithm}

% Array and table packages
\usepackage{array,multirow,booktabs,makecell}

% Graphics packages
\usepackage{graphicx,epsfig}

% Float and sizing packages
\usepackage{stfloats,float,calc}

% Text and formatting packages
\usepackage{textcomp,verbatim,url,cite,color,xcolor,cancel}
% Changes package for revision highlights
\let\comment\undefined
\usepackage[draft]{changes}
\definechangesauthor[name={cuiys}, color=red]{cuiys}
% 使用示例：\added[id=cuiys]{新增内容}，\deleted[id=cuiys]{删除内容}，\replaced[id=cuiys]{新}{旧}

% Layout packages
\usepackage{cuted}

\hyphenation{op-tical net-works semi-conduc-tor IEEE-Xplore}

% Environment customization
\makeatletter
\providecommand{\tabularnewline}{\\}
\floatstyle{ruled}
\newfloat{algorithm}{tbp}{loa}
\providecommand{\algorithmname}{Algorithm}
\floatname{algorithm}{\protect\algorithmname}
\makeatother

% Theorem environments
\theoremstyle{plain}
\newtheorem{theorem}{Theorem}
\newtheorem{assumption}{Assumption}
\newtheorem{lemma}{Lemma}
\newtheorem{lem}{\protect\lemmaname}
\newtheorem{thm}{\protect\theoremname}
\newtheorem{pro}{\protect\proposition}
\newtheorem{corollary}{Corollary}

\theoremstyle{definition}
\newtheorem{defn}{\protect\definitionname}
\newtheorem{myDef}{Definition}

% Custom command definitions
\providecommand{\definitionname}{Definition}
\providecommand{\lemmaname}{Lemma}
\providecommand{\theoremname}{Theorem}
\providecommand{\proposition}{Proposition}
\renewcommand{\algorithmicrequire}{\textbf{Input:}}
\renewcommand{\algorithmicensure}{\textbf{Output:}}

% Display settings
\allowdisplaybreaks[4]

\begin{document}

\title{Satellite TT\&C System Federated Fault Diagnosis:\\Architecture Design and Convergence Analysis}

\author{\IEEEauthorblockN{Yongsheng Cui\IEEEauthorrefmark{1}, Taiyi Chen\IEEEauthorrefmark{1}, Haoran Xie\IEEEauthorrefmark{2}, and Yafeng Zhan\IEEEauthorrefmark{1,3}} \\
\IEEEauthorblockA{\IEEEauthorrefmark{1}Department of Electronic Engineering, Tsinghua University, Beijing, China\\
\IEEEauthorrefmark{2}Electric Power Research Institute, China Southern Power Grid, Guangzhou, China\\
\IEEEauthorrefmark{3}Beijing National Research Center for Information Science and Technology, Tsinghua University, Beijing, China}
\thanks{This work was supported by the National Natural Science Foundation of China under Grant 62131012.}
\thanks{Y. Zhan is the corresponding author (e-mail: zhanyf@mail.tsinghua.edu.cn).}}

\maketitle

\begin{abstract}
Satellites operate in harsh space environments for extended periods, 
making fault diagnosis a critical strategy for extending their lifespan. 
Federated learning based fault diagnosis using telemetry data addresses critical challenges in satellite system reliability, 
providing solutions for data privacy across different telemetry, tracking and command (TT\&C) centers 
and improving fault diagnosis performance when individual TT\&C centers have limited training data. 
However, challenges arise in the deployment and application of federated learning methods
due to \textit{``Communication-limited''} and \textit{``Data Heterogeneity''} at some TT\&C centers, 
causing performance degradation and slow convergence.
To solve these problems, this paper proposes a streaming federated learning framework for satellite fault diagnosis. 
The framework incorporates an Incremental Learning Algorithm based on Dual-Domain Knowledge Distillation for efficient model updates, 
and a Federated Learning Algorithm based on Partial Forgetting Compensation for handling data heterogeneity. 
Under this framework, TT\&C centers with communication constraints only need to participate in a single round of federated learning 
for the system to obtain a well-performing and stable model. 
Theoretical convergence analysis is provided to characterize the algorithm performance. 
Extensive simulations show that the proposed algorithm outperforms existing fault diagnosis approaches for satellite systems.


\end{abstract}

\begin{IEEEkeywords}
  Communication constraints, convergence analysis, data heterogeneity, federated learning, knowledge distillation, satellite fault diagnosis.
\end{IEEEkeywords}

\section{INTRODUCTION}\label{sec:Introduction}
% Importance of satellite system monitoring
\IEEEPARstart{S}{atellite} fault diagnosis is critical for mission success in space operations, 
where maintenance opportunities are virtually nonexistent once deployed. 
The Telemetry, Tracking, and Command (TT\&C) system provides telemetry streams (electrical, thermal and status data) that 
enable continuous spacecraft health monitoring~\cite{ganesan2021fault}. 
Data-driven approaches can extract fault features from these telemetry streams~\cite{xiao2020deep}, 
yet their efficacy remains constrained in TT\&C centers by several factors: 
fault occurrences are infrequent, resulting in highly imbalanced datasets with scarce representative samples, 
and models trained on historical telemetry data often exhibit poor generalization capabilities when confronted with novel fault patterns emerging from diverse platforms 
or varying operational conditions~\cite{xie2021leo,bieber2023generic}.

% Feasibility of aggregating data across multiple TT&C centers
Aggregating fault data across multiple TT\&C centers would ideally address data scarcity and expand fault pattern coverage~\cite{von2018data}. 
By pooling telemetry data, each center could access more fault samples, diverse fault manifestations across platforms, 
and novel fault types absent in its local dataset—all essential for improving generalization to unseen faults in future operations. 
However, such direct data aggregation faces significant barriers in practice due to stringent data privacy requirements in aerospace applications. 
Satellite telemetry often contains sensitive information related to national security and privacy interests, 
rendering cross-center data sharing impractical.

% Introduction to federated learning
Federated learning~\cite{mcmahan2017communication} provides a promising privacy preserving alternative that reconciles these competing demands. 
By enabling collaborative model training without raw data exchange, 
federated learning preserves data sovereignty while allowing TT\&C centers to collectively benefit from shared fault pattern knowledge. 
In this paper, we investigate federated learning for satellite fault diagnosis, 
addressing two unique challenges specific to this domain:

\begin{figure}[t]
	\centering
	\includegraphics[width=\columnwidth]{fig/System Model/system model.pdf}
	\caption{System model.}
	\label{fig:system_model}
\end{figure}

% Problems addressed in this paper
\textit{Communication-limited}: A subset of TT\&C centers can participate in only a very small number of federated update rounds 
— in extreme cases only a single round — due to limited link availability, or operational constraints. 

\textit{Data Heterogeneity}: Telemetry data exhibits pronounced heterogeneity in two dimensions: 
(1) \textbf{feature-level}—identical faults manifest differently across centers due to platform and sensor diversity; 
(2) \textbf{label-level}—centers maintain disparate fault inventories with imbalanced distributions.

These challenges are illustrated in Fig.~\ref{fig:system_model}, where a communication-constrained TT\&C center joins at round 100 with new fault categories. 
Traditional federated learning then encounters: 
(1) \textbf{Model instability}—the model has difficulty retaining knowledge of old classes when learning new ones; 
(2) \textbf{Recency bias}—after the constrained center departs, the model gradually forgets its unique fault classes, as shown in Fig.~\ref{fig:problem-formulation}.

\begin{figure}[t]
	\centering
	\includegraphics[width=\columnwidth]{fig/Introduction/problem-formulation-acc.pdf}
  \caption{Problem formulation. $\alpha$ represents the degree of data heterogeneity.}
	\label{fig:problem-formulation}
\end{figure}

% Objective of this paper
The objective of this paper is to develop a model capable of accurately diagnosing faults across all TT\&C centers 
while minimizing the communication rounds needed for communication-constrained TT\&C centers. 
Related research on continual learning, federated learning with privacy concerns, and heterogeneous data handling is discussed in Section~\ref{sec:Related_Work}.

% Main contributions of this paper
In this study, we present a streaming federated learning framework for satellite fault diagnosis 
that ensures efficient communication and maintains data privacy. 
Under this framework, the TT\&C center with communication constraints only needs to 
participate in a \textbf{single round} of federated learning for the system to obtain a well-performing and stable model. 
% Summary of contributions
The key contributions of this work are outlined below:
\begin{itemize}
  \item 
  We provide a rigorous convergence analysis for the proposed algorithms in the communication-constrained satellite TT\&C setting. 
  Specifically, for the proposed algorithms under non-convex objectives, 
  we establish sublinear convergence to a neighborhood of a stationary point, 
  with a bias-controlled asymptotic error floor.
  \item 
  We design a streaming federated learning framework with two core algorithms: 
  (1) an Incremental Learning Algorithm based on Dual-Domain Knowledge Distillation for effective knowledge acquisition; 
  (2) a Federated Learning Algorithm based on Partial Forgetting Compensation to mitigate knowledge loss.
  \item 
  We perform extensive experiments, including comparisons with existing methods in the literature such as DTST and in-depth ablation studies. 
  Results show the framework's superior robustness under communication constraints and data heterogeneity.
\end{itemize}

% Paper organization
The remainder of this paper is structured as follows: 
Section \ref{sec:Related_Work} reviews related work on continual learning and federated learning. 
In Section \ref{sec:proposed_framework}, the proposed framework is detailed,
including the implementation procedures and the two core algorithms involved in this framework. 
Section \ref{sec:convergence-analysis} provides theoretical convergence analysis of the proposed algorithm.
Section \ref{sec:simulation-results} gives a series of case studies to demonstrate the effectiveness of the developed schemes. 
Section \ref{sec:conclusion} concludes the paper.
\input{chapter/Related Work.tex}
\section{STREAMING FEDERATED LEARNING FOR SATELLITE FAULT DIAGNOSIS}\label{sec:Proposed Streaming Federated Learning Framework}
% Framework overview and workflow diagram

% Detailed implementation stages
The streaming federated learning framework for satellite fault diagnosis 
is illustrated in Fig. \ref{fig:streaming-federated-learning-framework}. 
Whenever a new type of fault emerges at a communication-constrained TT\&C center 
or a new communication-constrained TT\&C center joins the system, the framework is re-initiated. The convergence time of federated learning is significantly shorter than the interval between new fault occurrences, enabling continuous execution of the framework in a streaming fashion. The framework operates through six distinct stages:

\begin{figure}[t]
	\centering
	\includegraphics[width=\columnwidth]{fig/Proposed Streaming Federated Learning Framework/Streaming federated learning framework.pdf}
	\caption{Streaming federated learning framework.}
	\label{fig:streaming-federated-learning-framework}
\end{figure}

\textit{Stage 1}: All unconstrained TT\&C centers first retrieve the global model from the central server 
and then conduct local training using their individual satellite fault datasets. 
After finishing local training, each center transmits its updated model parameters to the central server. 
Once the server receives parameters from all unconstrained TT\&C centers, 
it executes a model aggregation procedure to generate updated global model parameters, 
which are then redistributed to all participating centers. This process iterates until convergence criteria are met.

\textit{Stage 2}: The central server disseminates the converged global model (gray), 
containing the accumulated knowledge from all unconstrained TT\&C centers, to every communication-limited TT\&C center. 
This model serves as the foundational knowledge base for subsequent incremental learning processes in communication-constrained environments.

\textit{Stage 3}: Execute the Incremental Learning Algorithm based on Dual-Domain Knowledge Distillation 
(IL-DDKD) described in Section~\ref{sec:Proposed Streaming Federated Learning Framework}-\ref{section:3-1}. 
During this stage, the global model (gray) functions solely in inference mode, 
serving as a knowledge source to guide communication-constrained TT\&C centers in learning their respective partial models (blue). 
This knowledge transfer enables communication-constrained centers to acquire specialized capabilities 
for handling new fault types while preserving previously learned knowledge.

\textit{Stage 4}: Each communication-limited TT\&C center transmits its trained partial model (blue) to the central server for storage. 
The central server maintains a repository of partial models, where each model corresponds to a specific communication-constrained center. 
When a TT\&C center uploads a new partial model, the central server replaces any existing model from the same center, 
ensuring that only the most recent knowledge is retained for subsequent federated learning rounds.

\textit{Stage 5}: Execute the Federated Learning Algorithm based on Partial Forgetting Compensation 
(FL-PFC) detailed in Section~\ref{sec:Proposed Streaming Federated Learning Framework}-\ref{section:3-2}. 
During this stage, the stored partial models (blue) actively participate in the federated learning process 
by assisting unconstrained TT\&C centers with model training. 
Simultaneously, these partial models contribute to the central server's model parameter aggregation and global model updates. 
The global model undergoes continuous refinement and enhancement through this collaborative learning mechanism, 
transitioning from its initial state (gray) to an improved state (orange).

\textit{Stage 6}: The central server periodically synchronizes the enhanced global model (orange) with all communication-constrained TT\&C centers. 
This periodic update ensures that communication-constrained centers maintain access to the latest fault diagnosis capabilities, 
enabling them to benefit from the collective knowledge acquired through the federated learning process. 
The synchronization rate may be modified according to network conditions and fault detection requirements.

% Algorithm 1: Incremental Learning Algorithm Based on Dual-Domain Knowledge Distillation
\subsection{Incremental Learning Algorithm Based on Dual-Domain Knowledge Distillation}\label{section:3-1}
Building upon the Learning Without Forgetting approach~\cite{li2017learning} originally designed for task incremental learning~\cite{masana2022class}, 
we develop a dual-domain knowledge distillation~\cite{hinton2015distilling} based incremental learning technique tailored for class incremental challenges in satellite fault diagnosis. 
Our method employs a ``teacher-student'' paradigm, 
wherein the teacher model imparts its accumulated expertise to the student model, 
empowering the student to master new categories while retaining existing knowledge.
Due to the limited expressive power of label information, 
we obtain more distinctive feature representations from the intermediate layers of the teacher model.
The architectural diagram of this method is presented in Fig. \ref{fig:knowledge-distillation}.
Within this structure, the teacher model aligns with the original global model (gray) shown in Fig. \ref{fig:streaming-federated-learning-framework}, 
whereas the student model represents the partial model intended for training (blue) as illustrated in Fig. \ref{fig:streaming-federated-learning-framework}.
\begin{figure}[t]
	\centering
	\includegraphics[width=\columnwidth]{fig/Proposed Streaming Federated Learning Framework/Incremental learning algorithm based on dual-domain knowledge distillation.pdf}
	\caption{Incremental learning algorithm based on dual-domain knowledge distillation.}
	\label{fig:knowledge-distillation}
\end{figure}

Given the potential performance discrepancy between teacher and student models, 
students may struggle to efficiently absorb knowledge from teachers~\cite{mirzadeh2020improved}. 
To overcome this limitation, we employ module sharing between teacher and student architectures to enhance knowledge transfer efficiency. 
When dealing with new fault categories, additional nodes are incorporated into the network's output layer, with connection weights between these nodes and the previous layer randomly initialized. 
The loss function formulation during network training is expressed as follows:

\textbf{Supervised Loss}: This loss function is formulated to allow the partial model 
to efficiently learn fault data from the newly added TT\&C center. 
The cross entropy loss function is utilized for this objective:
\begin{equation}
	{\mathcal {L_C}}\left( {{\theta _{stu}}; {{\mathscr D}_{new}}} \right) =  - {{\bf{1}}_{{\mathcal Y_{new}}}} \log { \tilde{\mathcal Y}_{new}},
	\label{eq: LC}
\end{equation}
\noindent where ${\theta _{stu}}$ denotes the parameters of the student model,
${{\bf{1}}_{{\mathcal Y_{new}}}}$ represents the one-hot encoding of the ground truth label ${\mathcal Y_{new}}$,
${\tilde{\mathcal Y}_{new}}$ represents the model output after Softmax activation processing,
and ${{\mathscr D}_{new}}$ denotes the fault dataset from the new TT\&C center 
containing input features ${\mathcal X_{new}}$ and their corresponding true labels ${\mathcal Y_{new}}$.

\textbf{Dual domain distillation loss}: 
This loss component enables the partial model to continually assimilate fault knowledge from other TT\&C centers via the global model. 
Knowledge transfer is achieved by forwarding the attention maps produced by the intermediate layers of the partial model to the local model, 
directing the local model to concentrate on identical areas as the partial model, 
thus enabling data-domain distillation. 
Given that straightforward 1:1 alignment between feature representations of teacher and student models lacks theoretical grounding, 
we incorporate a regressor between them to enhance the alignment procedure \cite{romero2014fitnets}. 
Such architecture offers adaptability in the student model's training phase, 
preventing simple replication of the teacher model's representations by the student model.
The loss function for data-domain distillation ${\mathcal L_{f}}$ is defined as follows:
\begin{equation}
	{\mathcal L_{f}} = \frac{1}{2}{\left\| {u\left( {{{\theta _{tea}}}; {\mathcal X_{new}}} \right) - \gamma \left( {v\left( {\theta _{stu}; {\mathcal X_{new}}} \right); {W_r}} \right)} \right\|^2},
	\label{eq: LKD-1}
\end{equation}
\noindent where ${\theta _{tea}}$ represents the teacher model parameters, 
$\gamma $ denotes the regressor, ${W_r}$ refers to the parameters of the regressor, 
$u$ and $v$ represent the feature vectors generated by the teacher model and the student model, respectively. 

Meanwhile, this paper also performs distillation in the parameter domain by imposing constraints on key parameters 
to limit potential drastic changes during the learning process, 
thereby mitigating the interference of new knowledge with existing knowledge. 
The distillation loss function in the parameter domain ${\mathcal L_{par}}$ is defined as follows:
\begin{equation}
	{\mathcal L_{par}} = \frac{1}{2}\sum\limits_i {{\mathcal F_i}{{\left( {\theta _{stu}^i - \theta _{tea}^i} \right)}^2}},
	\label{eq: LKD-2}
\end{equation}
where $\mathcal F\left( \cdot \right)$ denotes the Fisher information matrix, and $i$ indicates the index of each parameter. 
The data-domain distillation loss in \eqref{eq: LKD-1} and the parameter-domain distillation loss in \eqref{eq: LKD-2} 
are combined as the loss function of dual-domain distillation:
\begin{equation}
	{\mathcal {L_{KD}}} = {\mathcal L_{f}}\left( {{\theta _{stu}};{\theta _{tea}};{W_r},{\mathcal X_{new}}} \right) + {\mathcal L_{par}}\left( {{\theta _{stu}};{\theta _{tea}}} \right).
\end{equation}

The supervised loss focuses on the performance of the algorithm on new fault classes, 
whereas the dual-domain distillation loss concentrates on the algorithm's performance regarding previous fault classes. 
The combination of these two components constitutes the overall training objective:
\begin{equation}
	{\mathcal L_{Dual}} = {\mathcal {L_C}}\left( {{\theta _{stu}};{{\mathscr D}_{new}}} \right) + {\lambda _o}{\mathcal {L_{KD}}}\left( {{\theta _{stu}};{\theta _{tea}};{W_r},{\mathcal X_{new}}} \right),
\end{equation}
where ${\lambda _o} > 0$ denotes the weighting parameter employed to balance the algorithm's performance between previous faults and new faults.

This paper employs an adaptive adjustment strategy for the weight coefficient ${\lambda _o}$.
Specifically, prior to training the student model, 
the teacher model is first duplicated and fine-tuned on the new TT\&C center's fault dataset ${\mathscr D}_{new}$ 
to obtain the baseline recognition accuracy $\mathcal{A}*$ on the new TT\&C center's faults. 
During the student model training process, a relatively large initial weight coefficient $\lambda_0$ 
is set to ensure minimal forgetting of previously learned faults, 
while establishing a performance threshold $\varTheta$ and a decay function $\mathcal{G}$. 
Subsequently, the recognition accuracy of the student model on the new TT\&C center's fault data is continuously monitored. 
If the student model's recognition accuracy on the new TT\&C center's fault data 
fails to reach the predefined performance threshold $(1-\varTheta)\mathcal{A}*$, 
the weight coefficient $\lambda_0$ is decayed according to $\lambda_0\rightarrow \mathcal{G}(\lambda_0)$, 
until the student model meets the expected performance requirements 
on the new TT\&C center's fault recognition task.

Algorithm \ref{Alg:IL-DDKD} presents the complete execution process of the IL-DDKD algorithm.
\begin{algorithm}[t]
	\caption{IL-DDKD Algorithm}
	\label{Alg:IL-DDKD}
	\begin{algorithmic}[1]
		\REQUIRE New TT\&C center fault data ${\mathcal X_{new}}$ and ground truth ${\mathcal Y_{new}}$, global model ${\theta _{tea}}$, local epochs $E$, local minibatch size $B$, accuracy drop margin $\Theta $, decay function $\mathcal G$, local learning rate $\eta_L $
		\ENSURE Partial model parameters $\theta _{{stu}}^ * $
		\STATE \textbf{Initialization}: Set the initial weight coefficient ${\lambda _o} = 100$, ${\mathcal A^ * } = 0$, ${W_r}$ to small random values.
		\STATE ${\mathcal A^ * } \leftarrow $Finetune$\left( {{{\mathscr D}_{new}},\eta_L ;{\theta _{tea}}} \right)$
		\WHILE {$\mathcal A < \left( {1 - \Theta } \right){\mathcal A^ * }$}
		\STATE $\mathcal B \leftarrow $ (Split ${{\mathscr D}_{new}}$ into batches of size $B$)
		\FOR {$e = 1, 2, 3,  ... , E$  }
		\FOR {batch $b \in \mathcal B$}
		\STATE ${\mathcal {L_C}} \leftarrow $CE$ ( {{{\bf{1}}_{{\mathcal Y_{new}}}},{\mathcal{\tilde Y}_{new}}} )$
		\STATE $teacher\_out \leftarrow u\left( {{\theta _{tea}},{\mathcal X_{new}}} \right)$
		\STATE $student\_out \leftarrow \gamma \left( {v\left( {{\theta _{stu}},{\mathcal X_{new}}} \right),{W_r}} \right)$
		\STATE ${L_{f}} \leftarrow$ MSE $\left( {teacher\_out, student\_out} \right)$
		\STATE ${\mathcal L_{par}} \leftarrow \mathcal F\left( {{\theta _{stu}},{\theta _{tea}}} \right)$
		\STATE ${\mathcal{L_{KD}}} = {\mathcal L_{f}} + {L_{par}}$
		\STATE ${\mathcal L_{Dual}} = {\mathcal {L_C}} + {\lambda _o}{\mathcal {L_{KD}}}$
		\STATE ${\theta _{stu}} \leftarrow {\theta _{stu}} - \eta_L \nabla {\mathcal L_{Dual}}$
		\ENDFOR
		\ENDFOR
		\STATE $\mathcal A \leftarrow $IL-DDKD$\left( {{{{\mathscr D}_{new}}};{\theta _{stu}}} \right)$
		\IF {$\mathcal A < \left( {1 - \Theta } \right){\mathcal A^ * }$}
		\STATE ${\lambda _o} \leftarrow \mathcal G\left( {{\lambda _o}} \right)$
		\ENDIF
		\ENDWHILE
		\STATE $\theta _{{{stu}}}^ *  \leftarrow {\theta _{stu}}$
		\RETURN $\theta _{{{stu}}}^ * $
	\end{algorithmic}
\end{algorithm}

% Algorithm 2: Federated Learning Algorithm Based on Partial Forgetting Compensation
\subsection{Federated Learning Algorithm Based on Partial Forgetting Compensation}\label{section:3-2}
\begin{figure}[t]
	\centering
	\includegraphics[width=\columnwidth]{fig/Proposed Streaming Federated Learning Framework/Federated learning algorithm based on local forgetting compensation.pdf}
	\caption{Federated learning algorithm based on partial forgetting compensation.}
	\label{fig:partial-forgetting-compensation}
\end{figure}

To enable the knowledge acquired at communication-constrained TT\&C centers 
to be utilized in subsequent federated learning processes of unconstrained TT\&C centers, 
this paper implements a virtual client as a proxy server for each communication-constrained TT\&C center. 
The proxy server stores the model trained by the communication-constrained TT\&C center 
and applies partial forgetting compensation to the entire system, 
ensuring that model parameter updates remain within the target range.
Considering the privacy leakage risks associated with model parameters or gradients~\cite{hitaj2017deep}, 
the proxy server provides white-box access exclusively to the highly trusted central server, 
while offering black-box access to individual TT\&C centers. 
The architecture of the federated learning algorithm based on partial forgetting compensation is illustrated in Fig. \ref{fig:partial-forgetting-compensation},
and the complete execution process is detailed in Algorithm \ref{Alg:FL-LFC}. 
The algorithm achieves model parameter updates through iterative interactions between the central server and the TT\&C centers. 
Each round of federated learning encompasses the following four steps:

\textbf{Global Model Download}: The central server chooses a subset of TT\&C centers for distributing the global model $\omega_t$ 
(where $t$ denotes the federated learning round) 
according to McMahan et al. \cite{mcmahan2017communication}, who found that selecting only a subset of clients for gradient descent 
in each federated learning round can yield a regularization effect comparable to Dropout. 
This paper employs a random selection strategy for TT\&C centers. 
However, if information regarding the effectiveness, reliability, and other relevant aspects of each TT\&C center becomes available, 
more sophisticated selection strategies can be developed based on this information to further enhance algorithm performance.

\textbf{Local Model Update}: The selected TT\&C centers utilize ${\omega _t}$ as the initial model and perform training 
on their respective local fault data ${{\mathscr D}_{i}}$ using the \textit{inheritance distillation method} 
described in Section~\ref{sec:Proposed Streaming Federated Learning Framework}-\ref{section:3-2}-\ref{section:3-2-1}, generating an updated local model $\omega _{t + 1}^i$.

\textbf{Local Model Upload}: Participating TT\&C centers transmit their updated local model parameters to the central server.

\textbf{Global Model Update}: The central server consolidates model parameters received from all participating TT\&C centers 
and refreshes the global model employing the \textit{Dynamic Weighting Method} elaborated in Section~\ref{sec:Proposed Streaming Federated Learning Framework}-\ref{section:3-2}-\ref{section:3-2-2}.

\begin{algorithm}[t]
	\caption{FL-PFC Algorithm}
	\label{Alg:FL-LFC}
	\begin{algorithmic}[1]
		\REQUIRE The $i$-th TT\&C center fault data ${{\mathscr D}_{i}}$, the number of TT\&C centers $K$, communication rounds $T$, local epochs $E$, fraction of TT\&C centers on each round $p$, model saved in proxy server $\theta _ \times ^ * $, global learning rate $\eta $
		\ENSURE Global model parameters ${\omega ^ * }$
		\FOR {$t = 1, 2, 3,  ... , T$}
		\STATE $m \leftarrow \max \left( {p \cdot K,1} \right)$
		\STATE ${S_t} \leftarrow $(randomly selected subset of $m$ TT\&C centers)
		\FOR {$i \in {S_t}$ \textbf{in parallel}}
		\STATE $\omega _{t+1}^i \leftarrow $\textbf{Heritage Distillation Method}(${{\mathscr D}_{i}}$; $\theta _ \times ^ * $; ${\eta}$; $E$)
		\ENDFOR
		\STATE ${\omega _{t + 1}} \leftarrow $\textbf{Dynamic Weighting Method}(${{\mathscr D}_{i}}$;${S_t}$; $\omega _{t}$; $\omega _{t+1}^i$)
		\ENDFOR
		\STATE ${\omega ^ * } \leftarrow {\omega _T}$
		\RETURN ${\omega ^ * }$
	\end{algorithmic}
\end{algorithm}

% Inheritance Distillation Method
\subsubsection{Inheritance Distillation Method}
\label{section:3-2-1}
This section maintains the teacher-student model paradigm. 
However, during this stage, the roles of teacher and student models are reversed: 
the teacher model signifies the model developed by the communication-limited TT\&C center,
whereas the student model transforms into the model developed by unconstrained TT\&C centers. 
As the entire mechanism embodies the progressive dissemination and preservation of knowledge, 
we designate it as ``\textbf{inheritance distillation}". 
Within this configuration, the teacher model aligns with the trained partial model (blue) illustrated in Fig.~\ref{fig:streaming-federated-learning-framework}, 
and the student model matches the original global model (gray) depicted in Fig.~\ref{fig:streaming-federated-learning-framework}.

During the inheritance distillation process, the teacher model is managed by the proxy server, 
which assists each unconstrained TT\&C center in processing local fault data and training the model in a black-box manner. 
The purpose is to separate the model from the knowledge representation, avoiding the exposure of sensitive information and achieving privacy protection. 
The algorithm's loss function is formulated as follows:

\textbf{Classification Loss}: This loss function is designed to enable the local model to effectively learn the fault data ${{\mathscr D}_i}$ from its respective TT\&C center. 
The classification loss $\mathcal L_C^i$ for TT\&C center $i$ is defined as follows:
\begin{equation}
	\label{eq: LCi}
	\mathcal L_C^i = {{\mathcal D}_{\rm{CE}}}\left( {{\mathcal P_i}\left( {\omega _{t,e}^i,{\mathcal X_i}} \right),{\mathcal Y_i}} \right),
\end{equation}
\noindent where $\omega _{t,e}^i$ represents the model parameters of TT\&C center $i$ at iteration $e$ of round $t$, 
$e$ denotes the local training iterations, 
$t$ represents the federated learning rounds, 
${\mathcal P_i}\left( {\omega _{t,e}^i, {\mathcal X_i}} \right)$ indicates the probability values predicted by the model $\omega _{t,e}^i$ for input data ${\mathcal X_i}$, 
${\mathcal Y_i}$ represents the ground truth labels, 
and ${{\mathcal D}_{\rm{CE}}}\left( { \cdot ,\cdot } \right)$ denotes the binary cross-entropy loss function.

\textbf{Distillation loss}: This loss function is designed to enable the local model to continuously acquire fault knowledge from the partial model. 
The knowledge transfer is accomplished by transmitting the attention maps generated from the intermediate layers of the partial model to the local model, 
guiding the local model to focus on the same regions as the partial model, 
thereby facilitating the learning of fault knowledge from the communication-constrained TT\&C center.
The distillation loss function $\mathcal L_{\mathcal KD}^i$ for TT\&C center $i$ is defined as follows:
\begin{equation}
	\label{eq: LKDi}
	\mathcal L_{\mathcal KD}^i = \frac{1}{2}{\left\| {{u_i}\left( {\theta _ \times ^ * ,{\mathcal  X_i}} \right) - {\gamma _i}\left( {{v_i}\left( {\omega _{t,e}^i,{\mathcal  X_i}} \right),W_r^i} \right)} \right\|^2},
\end{equation}
\noindent where $\theta _ \times ^ * $ represents the model saved in proxy server, 
$\gamma_i$ is the mapper, and $W^i_r$ are the model parameters of the regressor,
${u_i}$ and ${v_i}$ are the feature vectors output by the partial model and the local model, respectively. 

The integration of the classification loss in \eqref{eq: LCi} and the distillation loss in \eqref{eq: LKDi} constitutes the complete training objective:
\begin{equation}
	{\mathcal L^i} = \mathcal L_C^i\left( {\omega _{t,e}^i,{{\mathscr D}_i}} \right) + {\lambda ^i}\mathcal L_{\mathcal KD}^i\left( {\omega _{t,e}^i;\theta _ \times ^ * ;W_r^i,{\mathcal X_i}} \right),
\end{equation}
where $\lambda^i > 0$ is the weight coefficient 
used to balance the learning performance of the local model on its own TT\&C center's fault data 
and the communication-constrained TT\&C center's fault data.

Drawing on the approach from \cite{li2020federated}, this paper introduces a regularization term into the objective function 
to mitigate the impact of data heterogeneity among TT\&C centers. The final objective function of the algorithm is expressed as follows:
\begin{equation}
	{{\tilde {\mathcal  L^i}}} = \mathcal L_C^i + {\lambda ^i}\mathcal L_{\mathcal KD}^i + \frac{\mu }{2}{\left\| {\omega _{t,e}^i - {\omega _t}} \right\|^2},
\end{equation}
where ${\omega _t}$ denotes the global model parameters disseminated by the central server during the $t$-th round, 
and $\mu$ is the penalty constant that is adaptively adjusted \cite{li2020federated}. 
Specifically, this paper adopts a simple heuristic strategy: 
when the objective function value exhibits an upward trend, the penalty constant $\mu$ is increased to enhance the constraint effect; 
conversely, when the objective function value continues to decrease, 
the penalty constant $\mu$ is appropriately reduced to accelerate convergence. 
Algorithm \ref{Alg:Heritage Distillation Method} presents the complete execution process of the heritage distillation method.

\begin{algorithm}[t]
	\caption{Heritage Distillation Method}
	\label{Alg:Heritage Distillation Method}
	\begin{algorithmic}[1]
		\REQUIRE Fault data ${{\mathscr D}_{i}}$ from the $i$-th TT\&C center, the local model achieved by the $i$-th TT\&C center in the $t$-th round $\omega _{t + 1}^i$, global model for the $t$-th round ${\omega _t}$, the collection of TT\&C centers randomly chosen in the $t$-th round ${S_t}$
		\ENSURE The local model parameters update of the $i$-th TT\&C center $\omega _{t + 1}^i$
		\STATE \textbf{Initialization}: Set the weight coefficient ${\lambda ^i} = 1$, $W_r^i$ to small random values, $\omega _{t,0}^i = {\omega _t}$.
		\STATE $\mathcal B \leftarrow $ (Split ${{\mathscr D}_{i}}$ into batches of size $B$)
		\FOR {$e = 1, 2, 3,  ... , E$}
		\FOR {batch $b \in \mathcal B$}
		\STATE $\mathcal  L_C^i \leftarrow$ CE$\left( {{\mathcal P_i}\left( {\omega _{t,e}^i,{\mathcal X_i}} \right),{\mathcal Y_i}} \right)$
		\STATE ${\mathcal L_{\mathcal KD}^i} \leftarrow$ MSE$\left( {{u_i}\left( {\theta _ \times ^ * ,{\mathcal X_i}} \right),{\gamma _i}\left( {{v_i}\left( {\omega _{t,e}^i,{\mathcal X_i}} \right),W_r^i} \right)} \right)$
		\STATE ${ {\tilde { \mathcal L}^i}} = \mathcal L_C^i + {\lambda ^i} {\mathcal L_{\mathcal KD}^i} + \frac{\mu }{2}{\left\| {\omega _{t,e}^i - {\omega _t}} \right\|^2}$
		\STATE $\omega _{t,e + 1}^i \leftarrow \omega _{t,e}^i - {\eta}\nabla  {\tilde {\mathcal L}^i}$
		\ENDFOR
		\ENDFOR
		\STATE $\omega _{t + 1}^i \leftarrow \omega _{t,E}^i$
		\RETURN $\omega _{t + 1}^i$
	\end{algorithmic}
\end{algorithm}

% Dynamic Weighting Method
\subsubsection{Dynamic Weighting Method}
\label{section:3-2-2}
Following each client $k$'s local model update, 
the central server collects these model parameters and 
executes a weighted average to update the global model ${\omega _{t + 1}} \leftarrow \sum\nolimits_{k = 1}^K {\frac{{{n_k}}}{n}\omega _{t + 1}^k} $, 
which represents the standard approach recommended by traditional federated learning algorithms \cite{von2018data}. 
However, this conventional aggregation method demonstrates poor performance in our study, 
primarily because it assigns weights to each client based exclusively on data volume, 
without adequately considering the importance and individual characteristics of each client during global model aggregation. 
The distinctive challenge in our research stems from the data heterogeneity of data across TT\&C centers, 
which results in unsatisfactory global model performance.

\begin{algorithm}[t]
	\caption{Dynamic Weighting Method}
	\label{Alg:Dynamic Weighting Method}
	\begin{algorithmic}[1]
		\REQUIRE The $i$-th TT\&C center fault data ${{\mathscr D}_{i}}$, the local model obtained by the $i$-th TT\&C center in the $t$-th round $\omega _{t + 1}^i$, global model of the $t$-th round ${\omega _t}$, the set of TT\&C centers randomly selected in the $t$-th round ${S_t}$
		\ENSURE The global model parameters update of the central server ${\omega _{t + 1}}$
		\STATE $u_{t + 1}^i \leftarrow $Test (${{\mathscr D}_{i}}$, $\omega _{t + 1}^i$)
		\STATE $v_t^i \leftarrow $ Test (${{\mathscr D}_{i}}$, ${\omega _t}$)
		\STATE ${\gamma ^i} \leftarrow \% \; ({{e^{u_{t + 1}^i - v_t^i}}})  $
		\STATE ${\pi _{proxy}},{\gamma _t^{proxy}} \leftarrow \bar \pi ,\bar C_{\gamma}(t+1)^{-(1+\rho )}$
		\STATE ${\omega _{t + 1}} \leftarrow \sum\nolimits_{i \in {S_t} \cup proxy} {\phi \left[\kern-0.15em\left[ {{\pi _i}{\gamma ^i}}  \right]\kern-0.15em\right]\omega _{t + 1}^i} $
		\RETURN ${\omega _{t + 1}}$
	\end{algorithmic}
\end{algorithm}

This paper introduces an importance balancing factor $\gamma$ in the federated learning aggregation process 
to enable the global model to dynamically adjust the importance of knowledge contributions from different TT\&C centers throughout the learning process. 
The importance balancing factor $\gamma$ is defined as:
\begin{equation}
	{\gamma ^k} = \frac{{{e^{u_{t + 1}^k - v_t^k}}}}{{\sum\nolimits_{k \in {S_t}} {{e^{u_{t + 1}^k - v_t^k}}} }},
\end{equation}
\noindent where $u_{t + 1}^k$ and $v_t^k$ represent the classification accuracy of the local model at round $t + 1$ and the global model at round $t$ on the fault dataset ${{\mathscr D}_{k}}$ of the $k$-th TT\&C center, respectively. The global model parameter update on the central server is computed according to the following equation:
\begin{equation}
	{\omega _{t + 1}} \leftarrow \sum_{k \in {S_t} \cup proxy} {\phi \left[\kern-0.15em\left[ {{\pi _k}{\gamma ^k}} 
		\right]\kern-0.15em\right]\omega _{t + 1}^k} ,
\end{equation}
\noindent where ${\pi _k}$ denotes the proportion of fault data in each TT\&C center, 
and $\phi \left[\kern-0.15em\left[ \cdot \right]\kern-0.15em\right]$ represents the normalized mapping in the metric space.
Interestingly, in the model aggregation process, 
local models with poorer performance are intentionally assigned higher weights. 
An intuitive explanation is that models requiring more learning in previous rounds 
should receive greater attention in subsequent training iterations. 
This design effectively compensates for weak performers 
and thereby enhances the overall performance and robustness of the model.

Furthermore, throughout the entire federated learning process, 
unconstrained TT\&C centers gradually acquire 
the fault knowledge of the communication-constrained TT\&C center with assistance from the virtual client. 
As these centers continuously utilize TT\&C data from unconstrained centers for training, 
they can further improve their performance on local fault datasets. 
Consequently, the performance of local models trained by unconstrained TT\&C centers 
in each round will progressively surpass the model performance of the virtual client. 
At this stage, the virtual client essentially becomes a performance bottleneck. 
To address this issue, this paper incorporates an explicit polynomial decay schedule
\begin{equation}
	\gamma _t^{proxy}=\bar C_{\gamma}(t+1)^{-(1+\rho )}, \qquad \rho > 0,
\end{equation}
into the weight coefficient of the communication-constrained TT\&C center 
to mitigate the negative impact of the virtual client during model aggregation. This explicit schedule is also the form required by the convergence analysis to make the proxy residual term summable.
% \input{chapter/Convergence Analysis.tex}
\section{CONVERGENCE ANALYSIS}\label{sec:convergence-analysis}
\added[id=cuiys]{
In this section, we analyze the convergence behavior of the proposed federated learning algorithm under partial participation, dynamic weighting, and proxy-client tracking.
}
\subsection{Update Rule}
\added[id=cuiys]{
The goal is to minimize the global objective
$F(\omega) = \sum_{k=1}^K p_k F_k(\omega)$, where $p_k \ge 0$ and $\sum_{k=1}^K p_k = 1$.
At communication round $t$, a subset of clients $S_t$ is sampled; each selected client performs $E$ local SGD steps, and the resulting updates are aggregated.
}
\added[id=cuiys]{
To characterize the effects of partial participation and dynamic weighting, we decompose the practical global update into an ideal full-participation component and a deviation term.
Define
$g_t \triangleq \sum_{k=1}^K p_k \bar g_t^k$ as the ideal full-participation gradient, where $\bar g_t^k$ is the client-level average local gradient over $E$ steps.
Define
$d_t \triangleq \sum_{k \in S_t} \alpha_{k,t} \bar g_t^k$ as the aggregated gradient under dynamic weights $\alpha_{k,t}$.
The deviation is then
$b_t \triangleq d_t - g_t$.
}

\added[id=cuiys]{
The global model update can be written as
\begin{equation}
\omega_{t+1} = \omega_t - \eta_t E(g_t + b_t) + r_t \triangleq \omega_t + \Delta_t,
\end{equation}
where $\eta_t$ is the global learning rate and $r_t$ denotes the proxy-client residual term.
}
\subsection{Assumptions}
For the convergence analysis, we adopt the following standard assumptions.

\begin{assumption}[$L$-Smoothness]\label{ass:smoothness}
The global objective $F$ and all local objectives $F_k$ are differentiable and $L$-smooth.
\end{assumption}

\begin{assumption}[Bounded Local Variance]\label{ass:local-variance}
The stochastic gradients computed locally are unbiased and have bounded variance:
$\mathbb{E}\|g_{t,e}^k - \nabla F_k(\omega_{t,e}^k)\|^2 \le \sigma_l^2.$
\end{assumption}

\begin{assumption}[Bounded Data Heterogeneity]\label{ass:data-heterogeneity}
The divergence between local and global gradients is bounded such that
$\sum_{k=1}^K p_k \|\nabla F_k(\omega)\|^2 \le \kappa^2 + M\|\nabla F(\omega)\|^2,$
for some constants $\kappa^2 \ge 0$ and $M \ge 1$.
\end{assumption}

\begin{assumption}[Bias Decoupling]\label{ass:bias-decoupling}
The total deviation $b_t$ can be decomposed into a systematic bias and sampling variance:
$b_t = \mu_t + (b_t - \mu_t),$
where the expected systematic bias is bounded by $\|\mu_t\|^2 \le B_\mu^2$, and the sampling variance is bounded by
$\mathbb{E}_{S_t}\|b_t - \mu_t\|^2 \le B_V^2.$
\end{assumption}

\begin{assumption}[Residual Decay]\label{ass:residual-decay}
The proxy residual term decays sufficiently fast across rounds:
$\|r_t\|^2 \le \frac{D}{(t+1)^{2+2\rho}},$
where $D > 0$ and $\rho > 0$.
\end{assumption}

\subsection{Main Results}
\added[id=cuiys]{
Under the above assumptions, we now present the main convergence result for non-convex objectives.
}


\begin{theorem}[Convergence Bound]
\added[id=cuiys]{
Suppose Assumptions 1--5 hold. Let the learning rate be $\eta_t = \frac{\eta_0}{\sqrt{t+1}}$, with $\eta_0$ chosen such that $\eta_t E L \le c$ for a sufficiently small constant $c$. After $T$ communication rounds, the minimum expected squared gradient norm satisfies
}
\begin{align}
\min_{0 \le t < T} \mathbb{E}\|\nabla F(\omega_t)\|^2
&\le \mathcal{O}\!\left(\!\frac{\Delta_F}{E\sqrt{T}}\!\right)\!+\!\mathcal{O}\!\left(\!\frac{C_R}{E\sqrt{T}}\!\right)
\!+\!\mathcal{O}\!\left(\!B_\mu^2\!\right) \notag \\
&\!+\!\mathcal{O}\!\left(\!\frac{(\sigma_l^2\!+\!\kappa^2\!+\!B_V^2)E\ln T}{\sqrt{T}}\!\right),
\end{align}
\added[id=cuiys]{
where $\Delta_F = F(\omega_0) - F_{\inf}$, and $C_R < \infty$ is an absolute constant that bounds the infinite series induced by proxy residuals.
}
\end{theorem}

\added[id=cuiys]{
The complete proof is provided in the Appendix. Taking $T \to \infty$ yields the following asymptotic characterization.
}

\added[id=cuiys]{
\begin{corollary}[Asymptotic Error Floor]
Under the same conditions as Theorem 1, as $T \to \infty$, the algorithm converges to a neighborhood of a stationary point characterized by
\begin{equation}
\lim_{T \to \infty} \min_{0 \le t < T} \mathbb{E}\|\nabla F(\omega_t)\|^2 \le \mathcal{O}(B_\mu^2).
\end{equation}
\end{corollary}
}


\subsection{Discussion and Theoretical Insights}
\added[id=cuiys]{
Theorem 1 and Corollary 1 clarify how each error source contributes to optimization dynamics under dynamic partial participation.
By Assumption 5, the proxy residual term is summable, and its cumulative effect is bounded by a finite constant $C_R$.
After normalization by the effective progress term $\sum_t \eta_t = \Theta(\sqrt{T})$, this contribution decays as $\mathcal{O}(1/\sqrt{T})$.
}

\added[id=cuiys]{
Stochastic effects from local gradient noise ($\sigma_l^2$), data heterogeneity ($\kappa^2$), and client-sampling variance ($B_V^2$) scale with
$\sum_t \eta_t^2 = \Theta(\ln T)$.
Under the decaying stepsize schedule, these terms contribute
$\mathcal{O}\!\left(\frac{\ln T}{\sqrt{T}}\right)$
to the stationarity bound.
}

\added[id=cuiys]{
The systematic bias term $B_\mu^2$ remains non-vanishing and determines the asymptotic error floor.
Therefore, without an explicit debiasing mechanism, the method converges to a neighborhood of a stationary point rather than to an exact stationary point.
}
\section{SIMULATION RESULTS}\label{sec:simulation-results}
To verify the effectiveness of the streaming federated learning framework for satellite system fault diagnosis, 
this paper designs a simulation experiment with 1 central server node and 6 TT\&C center nodes.
The central server is a logical entity responsible for coordinating the federated learning process, including distributing global models, aggregating local model updates, and maintaining partial models from communication-constrained TT\&C centers. It can be implemented either as a dedicated physical server or as a virtual node in cloud environments.
Among the TT\&C centers, 5 centers have no communication constraints, 
while 1 TT\&C center has communication constraints (this configuration aligns with the scale of China's TT\&C centers).

\added[id=cuiys]{The fault data used in this study is generated through a Simulink-based simulation framework~\cite{liu2015fault} that is calibrated against real telemetry from real satellite system.
Specifically, we first perform comparative analysis between simulated signals and measured telemetry, and then tune the simulation parameters to align key characteristics, 
so that the generated data preserves realistic fault evolution patterns.}
During the data preprocessing stage, 
data augmentation techniques are integrated to broaden the training data distribution 
\added[id=cuiys]{and improve model robustness, with the explicit goal of mimicking realistic on-orbit operating conditions 
and sensing disturbances. Four domain-specific approaches are developed for satellite systems:}
(1) Random noise addition to simulate space environment interference;
(2) Random temporal shifts to mimic timing offsets in anomaly detection;
(3) Amplitude scaling to simulate variations in lighting and temperature conditions;
(4) Temporal masking to simulate telemetry data loss during operation.
\added[id=cuiys]{These augmentation techniques are combined into a unified pipeline to emulate practical scenarios 
that are difficult to fully capture in collected data, 
thereby enabling the model to better adapt to diverse temporal patterns and complex space environments encountered in real satellite operations.}

To emulate the Non-IID characteristics of fault data across TT\&C centers, 
we employ the Dirichlet distribution for data partitioning, where parameter $\alpha$ controls the distribution sharpness~\cite{hsu2019measuring}. 
Fig.~\ref{fig:alpha_comparison} visualizes how the dataset partitioning outcomes change as $\alpha$ varies.

\begin{figure*}[t]
    \centering
	\includegraphics[width=\textwidth]{fig/Simulation Results/alpha_comparison.pdf}
    \caption{Dataset partitioning results under different $\alpha$ values.}
    \label{fig:alpha_comparison}
\end{figure*}

Throughout the entire training procedure, 
the Adam optimizer \cite{2014Adam} is employed with Dropout regularization \cite{srivastava2014dropout} 
applied to fully connected layers to enhance model generalization capability. 
Xavier initialization \cite{glorot2010understanding} is utilized for initializing weights between newly added output layer nodes 
and their preceding layers. 
All experimental evaluations are conducted on a computing cluster 
comprising 8 NVIDIA GeForce RTX 3080 Ti GPUs, 
with the framework implemented using the PyTorch deep learning library.

% \added{In standard federated learning (corresponding to Stage 1), 
% 60\% of TT\&C centers are randomly selected for each training round, 
% with each center conducting 5 epochs of local model training on their respective datasets using a learning rate of 0.001.}

% \added{For model training at the communication-constrained TT\&C center (corresponding to Stage 3), 
% the training configuration comprises 100 rounds with a local batch size of 16 samples, 
% and a reduced learning rate of 0.0001 to accommodate the constrained communication environment.}

% \added{In federated learning among unconstrained TT\&C centers (corresponding to Stage 5), 
% the training protocol spans 300 global rounds with 60\% center participation per round. 
% Each participating center executes 5 local training epochs per federated round using a learning rate of 0.0001, 
% which is deliberately reduced from the initial training phase to ensure stable convergence 
% and optimal performance within practical epoch constraints.}

% \added{All simulations in this paper are conducted under an extreme yet representative scenario: 
% the communication-constrained TT\&C center simultaneously hosts fault categories 
% that are entirely disjoint from those present in other TT\&C centers. 
% This configuration mimics challenging real-world conditions where certain ground stations encounter unique fault patterns 
% not observed elsewhere in the network.}

\subsection{Performance Evaluation}
Under this configuration, the proposed framework is deployed 
to perform fault diagnosis for satellite systems across all participating TT\&C centers. 
The experimental results are illustrated in Fig. \ref{fig:performance_across_datasets}. 
Three distinct datasets are constructed for comprehensive evaluation: 
Dataset 1 contains fault data exclusively from unconstrained TT\&C centers; 
Dataset 2 includes fault data solely from the communication-constrained TT\&C center; 
Dataset 3 represents the aggregated dataset comprising all fault instances from every TT\&C center 
(containing nine fault classes in total, with the ninth class being exclusive to the constrained TT\&C center).

Fig. \ref{fig:performance_across_datasets} illustrates the model's behavior across multiple datasets through different training phases. 
The outcomes highlight two principal capacities of the proposed framework:
First, during the learning phase at communication-constrained TT\&C centers, 
the dual-domain knowledge distillation-based incremental learning approach empowers the model to efficiently gain internal expertise 
while predominantly maintaining formerly acquired external knowledge (as demonstrated by consistent performance on Dataset 1). 
Second, during the learning phase at unconstrained TT\&C centers, 
the forgetting compensation functionality offered by the proxy server permits the model to continually integrate fresh external knowledge 
while efficiently preserving internal knowledge (as confirmed by persistent performance on Dataset 2).  
When benchmarked against ``joint training" with all fault data centralized at a single TT\&C center—
representing the theoretical performance upper bound—the proposed method achieves competitive results.
An important distinction to note is that the training timeline depicted in Fig.~\ref{fig:performance_across_datasets} 
should not be confused with federated learning rounds; specifically, the communication-constrained TT\&C center 
engages in only a single federated learning round throughout the entire process.

\begin{figure}[t]
	\centering
	\includegraphics[width=\columnwidth]{fig/Simulation Results/main result.pdf}
	\caption{Performance across datasets.}
	\label{fig:performance_across_datasets}
\end{figure}

In practical engineering deployments, satellites continuously produce novel fault patterns over their operational lifetime. 
To examine how the proposed framework behaves as new faults emerge over time, 
we adopt the streaming federated learning scenario schematized in Fig.~\ref{fig:streaming_scenario}. 
The experiment unfolds in two phases.
In the first phase, a new fault class appears: the unconstrained TT\&C centers jointly hold seven fault classes, 
whereas the communication-constrained TT\&C center contains one fault class that is completely disjoint from those at the other centers. 
In the second phase, an additional fault class emerges; the unconstrained centers retain the same seven fault classes, while the constrained center now hosts two fault classes that remain unique relative to the unconstrained centers. 
This configuration mimics the continuous arrival of novel fault types during real satellite operations.
\begin{figure}[t]
    \centering
    \includegraphics[width=\columnwidth]{fig/Simulation Results/streaming_scenario.pdf}
	\caption{Streaming federated learning scenario: fault class distribution across TT\&C centers.}
    \label{fig:streaming_scenario}
\end{figure}

Fig. \ref{fig:streaming_federated_learning_performance} illustrates the performance evolution of the algorithm 
under streaming federated learning conditions. 
The proposed framework maintains robust performance continuity 
when encountering novel fault classes at communication-constrained TT\&C centers 
or when incorporating newly constrained centers into the federated learning process. 
Critically, the framework sustains consistent performance levels across iterative streaming sessions, 
exhibiting minimal degradation over successive iterations given sufficient model capacity. 
These empirical findings establish essential theoretical groundwork and provide compelling evidence 
for the framework's practical deployability in real-world satellite systems.

\begin{figure}[t]
	\centering
	\includegraphics[width=\columnwidth]{fig/Simulation Results/Scalability.pdf}
	\caption{Streaming federated learning performance.}
	\label{fig:streaming_federated_learning_performance}
\end{figure}

\subsection{Comparative Experiments}
We conduct a comprehensive performance evaluation by comparing the proposed approach with 
established fault diagnosis methodologies under identical experimental conditions. 
The comparative analysis focuses on three critical metrics: classification accuracy, 
communication efficiency (measured in federated learning rounds), and robustness to data heterogeneity. 
Fig. \ref{fig:comparison_with_sota_satellite_fault_diagnosis} presents the performance comparison between our proposed algorithm 
and contemporary satellite system fault diagnosis approaches, 
particularly the state-of-the-art Dual-aspect Time Series Transformer (DTST) method~\cite{xu2024dtst}, 
which represents a specialized space system fault diagnosis technique.
To ensure equitable comparison conditions, we acknowledge that in practical deployments, 
each TT\&C center typically encounters only a subset of all possible fault patterns. 
However, for the purpose of rigorous benchmarking, we adopt a modified experimental setup where all fault types are assigned to a single TT\&C center 
while maintaining extremely sparse sample representation for certain fault classes. 
This configuration enables meaningful performance assessment while preserving the essential characteristics of real-world data distribution challenges.

\begin{figure}[t]
	\centering
	\includegraphics[width=\columnwidth]{fig/Simulation Results/compare-1.pdf}
	\caption{Comparison with SOTA satellite fault diagnosis.}
	\label{fig:comparison_with_sota_satellite_fault_diagnosis}
\end{figure}

As demonstrated, the DTST algorithm, when trained on limited fault datasets, 
exhibits significant challenges in learning from data-scarce fault samples, consequently resulting in constrained recognition capabilities. 
Although such performance may satisfy conventional evaluation metrics (as many studies require $\ge 80$\% accuracy for fault diagnosis), 
it proves inadequate for practical deployment due to complete recognition failure for certain critical fault classes,
as empirically validated by the confusion matrix in Fig.~\ref{fig:confusion-matrix-dtst}. 
In contrast, our proposed methodology successfully amalgamates fault knowledge from all participating TT\&C centers, 
systematically overcoming the inherent limitations imposed by single-center data constraints regarding both fault diversity and sample sufficiency,
as illustrated in Fig.~\ref{fig:confusion-matrix-proposed}.

\begin{figure}[t]
	\centering
	\includegraphics[width=0.8\columnwidth]{fig/Simulation Results/confusion-matrix-20250328.pdf}
	\caption{Confusion matrix (DTST method).}
	\label{fig:confusion-matrix-dtst}
\end{figure}
\begin{figure}[t]
	\centering
	\includegraphics[width=0.8\columnwidth]{fig/Simulation Results/confusion-matrix-4-effectiveness.pdf}
	\caption{Confusion matrix (Proposed method).}
	\label{fig:confusion-matrix-proposed}
\end{figure}

Fig.~\ref{fig:comparison_with_federated_learning_methods} presents the performance comparison results between the proposed algorithm and the base Fed-Avg algorithm~\cite{mcmahan2017communication}. 
Experimental results demonstrate that compared to the traditional multi-iteration Fed-Avg approach, 
the proposed algorithm requires only a single federated learning iteration with the communication-constrained TT\&C center to train a well-performing and stable model. 
In contrast, when the standard Fed-Avg algorithm operates with only a single iteration involving the communication-constrained TT\&C center, 
its performance degrades significantly after that center ceases participation in federated learning due to the lack of forgetting compensation mechanisms. 
Additionally, due to the extreme heterogeneity of the data, the Fed-Avg algorithm exhibits highly unstable performance during the training process. 
However, when the proposed method is incorporated into the multi-iteration Fed-Avg algorithm, the stability of the overall approach can be significantly enhanced. 
This demonstrates that the suggested approach can also enhance the robustness of federated learning algorithms effectively.

\begin{figure}[t]
	\centering
	\includegraphics[width=\columnwidth]{fig/Simulation Results/compare-2.pdf}
	\caption{Comparison with federated learning methods.}
	\label{fig:comparison_with_federated_learning_methods}
\end{figure}


\subsection{Ablation Studies}
To rigorously validate the effectiveness of the proposed framework, 
comprehensive ablation studies were systematically conducted on two core components:
the incremental learning algorithm based on dual-domain knowledge distillation 
and the federated learning algorithm based on partial forgetting compensation. 
Table \ref{tab:ablation-study} presents the final accuracy results 
of the incremental learning algorithm ablation study based on dual-domain knowledge distillation in Stage 3. 
The experimental results demonstrate that fine-tuning methodologies suffer from catastrophic forgetting, 
causing significant performance degradation.
Feature extraction techniques preserve historical knowledge but limit the capacity to learn novel information, 
resulting in suboptimal fault discrimination performance. 
In contrast, our proposed incremental learning approach effectively acquires new knowledge 
while preserving existing knowledge integrity, 
achieving superior classification accuracy across all fault categories.

\begin{table}[t]
\caption{Ablation study for federated learning components}
\centering
\label{tab:ablation-study}
\resizebox{\columnwidth}{!}{
\begin{tabular}{c|c|c}
\hline
 & \textbf{Method} & \textbf{Accuracy} \\ \hline
\multirow{3}{*}{Stage 3} 
 & IL-DDKD & \textbf{0.978} \\ \cline{2-3}
 & Fine-Tuning & 0.933\\ \cline{2-3}
 & Feature Extraction & 0.894\\ \hline
\multirow{3}{*}{Stage 5} 
 & Heritage + Dynamic weight & \textbf{0.974}\\ \cline{2-3}
 & Heritage & 0.898\\\cline{2-3}
 & Dynamic weight & 0.958\\ \hline
\end{tabular}
}
\end{table}

Table \ref{tab:ablation-study} also provides the final accuracy outcomes 
of the federated learning algorithm ablation study grounded in partial forgetting compensation in Stage 5. 
The experimental findings indicate that the proxy server architecture introduced in this paper 
attains better performance maintenance compared to individual components.

\subsection{Robustness Analysis}
This section assesses the robustness of the suggested algorithm 
under different configurations of TT\&C center numbers and data heterogeneity degrees, 
thereby providing comprehensive insights into the framework's adaptability across diverse operational environments.

\begin{figure}[t]
	\centering
	\includegraphics[width=\columnwidth]{fig/Simulation Results/robustness-1.pdf}
	\caption{Performance stability across different $\alpha$ values.}
	\label{fig:algorithm_performance_vs_heterogeneity}

\end{figure}

Fig. \ref{fig:algorithm_performance_vs_heterogeneity} illustrates the algorithm's performance characteristics 
under varying degrees of data heterogeneity, as controlled by the parameter $\alpha$. 
Empirical results demonstrate that the proposed algorithm maintains robust performance stability for $\alpha \ge 1$, 
with only marginal performance degradation observed at $\alpha = 0.1$. 

\begin{figure}[t]
	\centering
	\includegraphics[width=\columnwidth]{fig/Simulation Results/robustness-2.pdf}
	\caption{Performance stability across different TT\&C centers}
	\label{fig:algorithm_performance_vs_client_count}
\end{figure}

Fig. \ref{fig:algorithm_performance_vs_client_count} illustrates the algorithm's performance characteristics 
under varying numbers of TT\&C centers, thereby evaluating its scalability properties. 
Experimental results demonstrate that the proposed algorithm exhibits minimal performance variation 
across different client configurations, indicating robust insensitivity to client count fluctuations. 
This superior scalability characteristic constitutes a fundamental advantage for large-scale deployment scenarios 
and provides essential technical foundation for future engineering implementations in diverse satellite network architectures.
\section{CONCLUSION}\label{sec:conclusion}
This paper proposes a streaming federated learning framework for satellite system fault diagnosis with efficient communication and privacy protection. 
To address the challenges of "how to learn" and "how to fuse" for TT\&C centers with communication constraints, 
two core algorithms are designed:
(1) an Incremental Learning Algorithm based on Dual-Domain Knowledge Distillation for effective knowledge acquisition;
(2) a Federated Learning Algorithm based on Partial Forgetting Compensation to mitigate knowledge loss.
Theoretical analysis shows that, the proposed algorithms converge to a neighborhood of a stationary point with an asymptotic error floor.
Experimental results demonstrate that the framework achieves near-centralized performance with only one federated round, 
maintains stable online learning over time, and significantly outperforms existing approaches in terms of fault diagnosis accuracy. 
This work establishes a solid foundation for practical deployment of satellite fault diagnosis systems 
and demonstrates the significant potential of federated learning in satellite fault diagnosis applications.
% \appendix
\section{Proof Details}\label{app:B}
\subsection*{Lemma~\ref{lem:1} (Client Drift Bound)}

\noindent\textbf{Step 1 --- Drift recursion.}
\(\omega_{t,e}^k - \omega_t = -\eta_t\sum_{i=0}^{e-1}g_{t,i}^k\).
Cauchy--Schwarz:
\(\|\omega_{t,e}^k - \omega_t\|^2 \le \eta_t^2 e\sum_{i=0}^{e-1}\|g_{t,i}^k\|^2\).
Taking \(\mathbb E_t\) and weighted summation:
\(D_e \le \eta_t^2 e\sum_{i=0}^{e-1}\sum_k p_k\,\mathbb E_t\|g_{t,i}^k\|^2\).

\noindent\textbf{Step 2a --- Noise decomposition.}
\(g_{t,i}^k = \nabla F_k(\omega_{t,i}^k) + \varepsilon_{t,i}^k\)
(\(\mathbb E_t[\varepsilon] = 0\), \(\mathbb E_t\|\varepsilon\|^2 \le \sigma_l^2\)).
By \(\|a+b\|^2 \le 2\|a\|^2 + 2\|b\|^2\):
\(\mathbb E_t\|g_{t,i}^k\|^2 \le 2\sigma_l^2 + 2\|\nabla F_k(\omega_{t,i}^k)\|^2\).

\noindent\textbf{Step 2b --- Pullback.}
\(\nabla F_k(\omega_{t,i}^k) = \nabla F_k(\omega_t) + (\nabla F_k(\omega_{t,i}^k) - \nabla F_k(\omega_t))\).
\(L\)-smoothness:
\(\|\nabla F_k(\omega_{t,i}^k)\|^2 \le 2\|\nabla F_k(\omega_t)\|^2 + 2L^2\|\omega_{t,i}^k - \omega_t\|^2\).

\noindent\textbf{Step 2c --- Combine.}
\(\mathbb E_t\|g_{t,i}^k\|^2 \le 2\sigma_l^2 + 4\|\nabla F_k(\omega_t)\|^2 + 4L^2\,\mathbb E_t\|\omega_{t,i}^k - \omega_t\|^2\).

\noindent\textbf{Step 2d --- Weighted summation.}
\(\sum_k p_k\,\mathbb E_t\|g_{t,i}^k\|^2
 \le 2\sigma_l^2 + 4\sum_k p_k\|\nabla F_k(\omega_t)\|^2 + 4L^2 D_i\).
Using the heterogeneity assumption~\ref{asmp:2}:
\(\le A_t + 4L^2 D_i,\; A_t \triangleq 2\sigma_l^2 + 4\kappa^2 + 4M\|\nabla F(\omega_t)\|^2\).

\noindent\textbf{Steps 3--5 --- Solve via prefix maximum.}
Substituting yields \(D_e \le \eta_t^2 e^2 A_t + 4L^2\eta_t^2 e\sum_{i=0}^{e-1}D_i\).
Defining \(D_{\max,e} \triangleq \max_{0\le j\le e} D_j\):
\(D_{\max,e} \le \eta_t^2 e^2 A_t + 4L^2\eta_t^2 e^2 D_{\max,e}\).
Under \(8L^2\eta_t^2 E^2 \le 1\), we have \(4L^2\eta_t^2 e^2 \le 1/2\),
so \(D_{\max,e} \le 2\eta_t^2 e^2 A_t\). \(\square\)

\subsection*{Lemma~\ref{lem:2} (Gradient Bias and Second Moment)}

\noindent\textbf{Bias term \(\delta_t\):}
\(\delta_t = \frac{1}{E}\sum_{e,k} p_k(\nabla F_k(\omega_{t,e}^k) - \nabla F_k(\omega_t))\).
By Jensen's inequality and \(L\)-smoothness:
\(\mathbb E_t\|\delta_t\|^2 \le \frac{L^2}{E}\sum_{e=0}^{E-1} D_e
 \le L^2 D_{\max,E} \le 2L^2\eta_t^2 E^2 A_t\).\;
This establishes Eq.\,\eqref{eq:30}.

\noindent\textbf{Second moment \(\mathbb E_t\|g_t\|^2\):}
\(\mathbb E_t\|g_t\|^2 \le A_t + 4L^2 D_{\max,E} \le (1+8L^2\eta_t^2 E^2)A_t\).
Substituting \(A_t = 2\sigma_l^2 + 4\kappa^2 + 4M\|\nabla F\|^2\):
\begin{align*}
  (1+8L^2\eta_t^2 E^2)A_t
  &= 2(1+8L^2\eta_t^2 E^2)\sigma_l^2
     + 4(1+8L^2\eta_t^2 E^2)\kappa^2 \notag\\
  &\quad + 4M(1+8L^2\eta_t^2 E^2)\,\|\nabla F\|^2.
\end{align*}
Using \(\eta_t \le \bar\eta\) and the definition of \(c_g\) (Eq.\,\eqref{eq:13}):
\(\frac{1}{2}c_g \ge 2(1+8L^2\bar\eta^2 E^2)\),
\(c_g \ge 4(1+8L^2\bar\eta^2 E^2)\), and
\(c_g \ge 4M(1+8L^2\bar\eta^2 E^2)\).
The three terms are bounded by \(\frac{c_g}{2}\sigma_l^2\), \(c_g\kappa^2\), and \(c_g\|\nabla F\|^2\),
summing to \(\le c_g(\sigma_l^2 + \kappa^2 + \|\nabla F\|^2)\). \(\square\)

\subsection*{Lemma~\ref{lem:3} (Single-Step Descent)}

\noindent\textbf{Step 1 --- Inner product.}
\(\mathbb E_t\langle\nabla F, \Delta_t\rangle
 = -\eta_t E\langle\nabla F, \mathbb E_t[g_t]\rangle
   - \eta_t E\langle\nabla F, m_t\rangle
   + \mathbb E_t\langle\nabla F, r_t\rangle\).
Using \(\mathbb E_t[g_t] = \nabla F + \delta_t\):
\begin{align*}
  \text{(A)}\;& -\eta_t E\langle\nabla F, \nabla F + \delta_t\rangle \notag\\
              &\qquad = -\eta_t E\|\nabla F\|^2 - \eta_t E\langle\nabla F, \delta_t\rangle\\
  \text{(B)}\;& -\eta_t E\langle\nabla F, m_t\rangle,\qquad
  \text{(C)}\; \mathbb E_t\langle\nabla F, r_t\rangle.
\end{align*}
Applying Young's inequality (with \(\alpha = 1/2\) for (A), \(\alpha = 1/2\) for (B), \(\alpha = \eta_t E/4\) for (C)):
\begin{align*}
  \text{(A)} &\le -\eta_t E\|\nabla F\|^2 + \tfrac{\eta_t E}{4}\|\nabla F\|^2 \notag\\
             &\qquad + \eta_t E\|\delta_t\|^2
              = -\tfrac{3\eta_t E}{4}\|\nabla F\|^2 + \eta_t E\|\delta_t\|^2,\\
  \text{(B)} &\le \tfrac{\eta_t E}{4}\|\nabla F\|^2 + \eta_t E\|m_t\|^2,\\
  \text{(C)} &\le \tfrac{\eta_t E}{8}\|\nabla F\|^2 \notag\\
            &\qquad + \tfrac{2}{\eta_t E}\,\mathbb E_t\|r_t\|^2.
\end{align*}
Summing:
\begin{align*}
  \mathbb E_t\langle\nabla F, \Delta_t\rangle
  &\le -\tfrac{3}{8}\eta_t E\|\nabla F\|^2 \quad + \eta_t E\|\delta_t\|^2 \notag\\
  &\quad + \eta_t E\!\cdot\! B_\mu^2
     + \tfrac{2}{\eta_t E}\,\mathbb E_t\|r_t\|^2.
\end{align*}

\noindent\textbf{Step 2 --- Quadratic term.}
\(\|\Delta_t\|^2 \le 3\eta_t^2 E^2\|g_t\|^2 + 3\eta_t^2 E^2\|b_t\|^2 + 3\|r_t\|^2\).
Multiplying by \(L/2\), taking \(\mathbb E_t\), and substituting bounds:
\begin{align*}
  \frac{L}{2}\mathbb E_t\|\Delta_t\|^2
  &\le \frac{3L}{2}\eta_t^2 E^2 c_g(\sigma_l^2 + \kappa^2 + \|\nabla F\|^2) \\
  &\quad + \frac{3L}{2}\eta_t^2 E^2(B_\mu^2 + B_V^2)
     + \frac{3L}{2}\,\mathbb E_t\|r_t\|^2.
\end{align*}

\noindent\textbf{Step 3 --- Merging.}
Substituting \(\eta_t E\|\delta_t\|^2 \le 8L^2\eta_t^3 E^3(\sigma_l^2 + \kappa^2 + M\|\nabla F\|^2)\) (from Eq.\,\eqref{eq:30}):
\begin{align*}
  \mathbb E_t[F(\omega_{t+1})]
  &\le F(\omega_t)
     - \tfrac{3}{8}\eta_t E\|\nabla F\|^2
     + \tfrac{3L}{2}\eta_t^2 E^2 c_g\|\nabla F\|^2 \notag\\
  &\quad + 8L^2\eta_t^3 E^3 M\|\nabla F\|^2 \notag\\
  &\quad + \bigl(\tfrac{3L}{2}c_g\eta_t^2 E^2 + 8L^2\eta_t^3 E^3\bigr)(\sigma_l^2 + \kappa^2) \\
  &\quad + \tfrac{3L}{2}\eta_t^2 E^2 B_V^2
     + \eta_t E B_\mu^2
     + \tfrac{3L}{2}\eta_t^2 E^2 B_\mu^2 \notag\\
  &\quad + \bigl(\tfrac{2}{\eta_t E} + \tfrac{3L}{2}\bigr)\,\mathbb E_t\|r_t\|^2.
\end{align*}

\noindent\textbf{Step 4 --- Coefficient absorption.}
The coefficient of \(\|\nabla F\|^2\) is
\(-\frac{3}{8}\eta_t E + \frac{3L}{2}\eta_t^2 E^2 c_g + 8L^2\eta_t^3 E^3 M
 = -\frac{3}{8}\eta_t E + \eta_t E \cdot \eta_t E L \cdot (\frac{3}{2}c_g + 8L\eta_t E M)\).
Under \(\eta_t E L \le c_0 \le \frac{1}{16(1+c_g)}\) and \(c_g \ge 4M\),
one can verify \(\frac{3}{2}c_g + 8L\eta_t E M \le \frac{7}{2}c_g\),
so the coefficient satisfies
\(\le -\frac{3}{8}\eta_t E + \eta_t E \cdot c_0 \cdot \frac{7}{2}c_g
   \le -\frac{\eta_t E}{8}\)
(since \(\frac{7}{2}c_0 c_g < \frac{1}{4}\)).

\noindent\textbf{Step 5 --- Constants \(c_1, c_2\).}
Using \(8L^2\eta_t^3 E^3 \le 8L c_0 \eta_t^2 E^2\) to reduce order:
\((\frac{3L}{2}c_g + 8L c_0)\eta_t^2 E^2(\sigma_l^2 + \kappa^2)
   + \frac{3L}{2}\eta_t^2 E^2 B_V^2
 \le c_1\eta_t^2 E^2(\sigma_l^2 + \kappa^2 + B_V^2)\),
with \(c_1 = \frac{3L}{2}c_g + 8L c_0 + 3L\).
The \(B_\mu^2\) term:
\(\eta_t E + \frac{3L}{2}\eta_t^2 E^2
 = \eta_t E + \frac{3}{2}(\eta_t E L)\eta_t E
 \le (1 + \frac{3}{2}c_0)\eta_t E\),
yielding \(c_2 = 2 + 3c_0\). \(\square\)

% =================================================================
\subsection{Complete Proofs of Theorems}
% =================================================================

\subsection*{Theorem~\ref{thm:1} (Surrogate Objective Convergence)}

\noindent\textbf{Step 1 (Constant setup).}
From Proposition~\ref{prop:2}: \(\|m_t\|^2 \le B_\mu^2 = 4G^2\),
\(\mathbb E_t\|\zeta_t\|^2 \le B_V^2 = 4G^2\).
From Proposition~\ref{prop:3}: \(\|r_t\|^2 \le D/(t+1)^{2+2\rho}\).

\noindent\textbf{Step 2 (Telescoping sum).}
Take full expectation of Eq.\,\eqref{eq:33} and sum.
Let \(S_T \triangleq \sum_{t=0}^{T-1}\eta_t\).
The five terms accumulate as:
\begin{itemize}[leftmargin=*,nosep]
  \item \(\mathbb E_t[F(\omega_{t+1})] \le \mathbb E_t[F(\omega_t)]
        \;\Longrightarrow\; \mathbb E[F(\omega_T)] \le \mathbb E[F(\omega_0)]\)
        (telescoping cancellation)
  \item \(-\dfrac{\eta_t E}{8}\|\nabla F(\omega_t)\|^2
        \;\Longrightarrow\; -\dfrac{E}{8}\sum_{t=0}^{T-1}\eta_t\,\mathbb E\|\nabla F(\omega_t)\|^2\)
  \item \(c_1\eta_t^2 E^2(\sigma_l^2 + \kappa^2 + B_V^2)
        \;\Longrightarrow\; c_1 E^2(\sigma_l^2 + \kappa^2 + 4G^2)\sum_{t=0}^{T-1}\eta_t^2\)
  \item \(c_2\eta_t E \cdot 4G^2
        \;\Longrightarrow\; 4c_2 G^2 E \cdot S_T\)
  \item \(\bigl(\dfrac{8}{\eta_t E} + \dfrac{3L}{2}\bigr)\mathbb E_t\|r_t\|^2
        \;\Longrightarrow\; \sum_{t=0}^{T-1}\bigl(\dfrac{8}{\eta_t E} + \dfrac{3L}{2}\bigr)\mathbb E\|r_t\|^2\)
\end{itemize}

\noindent\textbf{Step 3 (Rearranging).}
From \(F(\omega_T) \ge F_{\inf}\):
\(\frac{E}{8}\sum\eta_t\,\mathbb E\|\nabla F\|^2
 \le \Delta_F + c_1 E^2(\sigma_l^2 + \kappa^2 + 4G^2)\sum\eta_t^2
    + 4c_2 G^2 E \cdot S_T + \mathcal R_\infty\).

\noindent\textbf{Steps 4--5 (Extract minimum \& series estimates).}
\(\min_{0\le t<T} \le \frac{8}{E S_T}[\Delta_F + c_1 E^2(\cdots)\sum\eta_t^2 + 4c_2 G^2 E \cdot S_T + \mathcal R_\infty]\).
With \(\eta_t = \eta_0/\sqrt{t+1}\):
\(S_T \ge \eta_0\sqrt{T}\;(T\ge 4)\), \(\sum\eta_t^2 \le \eta_0^2(1 + \ln T)\).

\noindent\textbf{Step 6 (Four-term substitution).}

\(\displaystyle\frac{8}{E S_T}\Delta_F \le \frac{8\Delta_F}{E\eta_0\sqrt{T}}\)

\(\displaystyle\frac{8}{E S_T}\,c_1 E^2(\cdots)\sum\eta_t^2
           \le \frac{8c_1 E\eta_0(1+\ln T)}{\sqrt{T}}(\sigma_l^2 + \kappa^2 + 4G^2)\)

\(\displaystyle\frac{8}{E S_T}\cdot 4c_2 G^2 E \cdot S_T = 32 c_2 G^2 E\)
           (\textbf{error floor}, \(S_T\) cancels)
           
\(\displaystyle\frac{8}{E S_T}\mathcal R_\infty \le \frac{8\mathcal R_\infty}{E\eta_0\sqrt{T}}\)


\noindent\textbf{Step 7 (Residual convergence verification).}
\(\mathcal R_\infty
 \le D\Bigl[\frac{8}{E\eta_0}\sum(t+1)^{-3/2-2\rho}
           + \frac{3L}{2}\sum(t+1)^{-2-2\rho}\Bigr] < \infty\)
(since \(\rho > 0\)). \(\square\)

\subsection*{Theorem~\ref{thm:2} (Adaptive Objective Convergence)}

\noindent\textbf{Step 1.}
Proposition~\ref{prop:5} guarantees that Eq.\,\eqref{eq:33} holds for \(\widetilde F_t\).

\noindent\textbf{Step 2.}
From \(\widetilde F_{t+1}(\omega_{t+1}) \le \widetilde F_t(\omega_{t+1}) + \varepsilon_t\) (Assumption~\ref{asmp:4}),
taking full expectation:
\begin{align*}
  \mathbb E[\widetilde F_{t+1}]
  &\le \mathbb E[\widetilde F_t]
     - \tfrac{\eta_t E}{8}\,\mathbb E\|\nabla\widetilde F_t\|^2
     + c_1\eta_t^2 E^2(\cdots) \\
  &\quad + 4c_2 G^2\eta_t E
     + \bigl(\tfrac{8}{\eta_t E} + \tfrac{3L}{2}\bigr)\mathbb E\|r_t\|^2
     + \varepsilon_t.
\end{align*}

\noindent\textbf{Step 3 (Telescoping sum).}
Summing from \(t = 0\) to \(T-1\):
\(\frac{E}{8}\sum\eta_t\,\mathbb E\|\nabla\widetilde F_t\|^2
 \le \widetilde\Delta_0 + \sum\varepsilon_t
    + c_1 E^2(\cdots)\sum\eta_t^2
    + 4c_2 G^2 E \cdot S_T + \mathcal R_\infty\).
Dividing by \(S_T\) and taking the minimum:
\(\frac{8\sum\varepsilon_t}{E S_T} \le \frac{8\mathcal E_\infty}{E\eta_0\sqrt{T}}\);
the error floor \(32 c_2 G^2 E\) remains unchanged. \(\square\)

%%%%%%%%%%%%%%%%%%%%%%%%%%%%%%%%%%%%%%%%%%%%%%%%%%%%%%%%%%%%%%

\bibliographystyle{IEEEtran}
\bibliography{bibfile}

%%%%%%%%%%%%%%%%%%%%%%%%%%%%%%%%%%%%%%%%%%%%%%%%%%%%%%%%%%%%%%

\begin{IEEEbiography}[{\includegraphics[width=1in,height=1.25in,clip,keepaspectratio]{fig/picture/Cuiyongsheng.jpg}}]{Yongsheng Cui}
	(Graduate Student Member, IEEE)
	is currently pursuing the Ph.D. degree in the Department of Electronic Engineering, Tsinghua University, Beijing, China.
	His research interests include deep space communications, satellite constellation design,
	machine learning and satellite TT\&C systems.
\end{IEEEbiography}

\begin{IEEEbiography}[{\includegraphics[width=1in,height=1.25in,clip,keepaspectratio]{fig/picture/Chentaiyi.jpg}}]{Taiyi Chen}
	(Student Member, IEEE)
	received the Master degree in telecommunication engineering from the Department of Electronic Engineering,
	Tsinghua University, Beijing, P.R. China.
	His research interests include Wireless Communication Systems for satellite constellations,
	channel coding, and deep space communications.
\end{IEEEbiography}

\begin{IEEEbiography}[{\includegraphics[width=1in,height=1.25in,clip,keepaspectratio]{fig/picture/Haoran Xie-eps-converted-to.pdf}}]{Haoran Xie}
	(Member, IEEE)
	is a researcher at the Electric Power Research Institute, China Southern Power Grid.
	His research interests include application of optimization, machine learning,
	and statistical theory.
	He has served as the Chair for IEEE IWCMC 2022/2023/2024 Symposium and IEEE VTC 2022-Fall Symposium.
\end{IEEEbiography}

\begin{IEEEbiography}[{\includegraphics[width=1in,height=1.25in,clip,keepaspectratio]{fig/picture/YafengZhan-eps-converted-to.pdf}}]{Yafeng Zhan}
	(Member, IEEE)
	received B.S.E.E. and Ph.D.E.E. degrees from the Department of Electronic Engineering, Tsinghua University, Beijing, China,
	in 1999 and 2004, respectively.
	He is currently a Professor with the Beijing National Research Center for Information Science and Technology,
	Tsinghua University. His current research interests include satellite TT\&C systems,
	communication signal processing and deep space communications.
	He serves as an editor of \textit{Chinese Space Science and Technology} and the \textit{Journal of Deep Space Exploration}.
\end{IEEEbiography}

\section*{APPENDIX}

This appendix provides proofs for the convergence results in the main text.

\subsection{Auxiliary Lemmas}

\begin{lemma}[Bounding the Client Drift]
Under Assumptions 1--3, the expected distance between the local model at step $e$ and the global model at round $t$, denoted as
\begin{equation}
D_e = \sum_{k=1}^K p_k \mathbb{E}_t\|\omega_{t,e}^k - \omega_t\|^2,
\end{equation}
is bounded by:
\begin{equation}
D_e \le \mathcal{O}(\eta_t^2 E^2) \left( C_\sigma + \|\nabla F(\omega_t)\|^2 \right),
\end{equation}
where $C_\sigma \triangleq 2\sigma_l^2 + 4\kappa^2$, and $\mathbb{E}_t[\cdot]$ denotes the expectation conditional on $\omega_t$. Consequently, the bias of the ideal gradient, $\delta_t = \mathbb{E}_t[g_t] - \nabla F(\omega_t)$, satisfies
\begin{equation}
\mathbb{E}_t\|\delta_t\|^2 \le \mathcal{O}(L^2\eta_t^2 E^2)(C_\sigma + 4M\|\nabla F(\omega_t)\|^2).
\end{equation}
\end{lemma}

\begin{proof}
From the local update rule $\omega_{t,e}^k=\omega_t-\eta_t\sum_{i=0}^{e-1}g_{t,i}^k$ and Cauchy--Schwarz,
\begin{equation}
\|\omega_{t,e}^k - \omega_t\|^2 \le \eta_t^2 e \sum_{i=0}^{e-1} \|g_{t,i}^k\|^2.
\end{equation}
Taking conditional expectation and weighted averaging over clients gives
\begin{equation}
\label{eq:bound_drift}
\sum_{k=1}^K p_k \mathbb{E}_t\|\omega_{t,e}^k - \omega_t\|^2
\le \eta_t^2 e \sum_{i=0}^{e-1} \sum_{k=1}^K p_k \mathbb{E}_t\|g_{t,i}^k\|^2.
\end{equation}
Using variance decomposition and $\|a+b\|^2 \le 2\|a\|^2+2\|b\|^2$,
\begin{equation}
\mathbb{E}_t\|g_{t,i}^k\|^2 \le 2\sigma_l^2 + 2\mathbb{E}_t\|\nabla F_k(\omega_{t,i}^k)\|^2.
\end{equation}
Adding and subtracting $\nabla F_k(\omega_t)$ and applying Assumption \ref{ass:smoothness},
\begin{equation}
\mathbb{E}_t\|g_{t,i}^k\|^2 \le 2\sigma_l^2 + 4\|\nabla F_k(\omega_t)\|^2 + 4L^2\mathbb{E}_t\|\omega_{t,i}^k - \omega_t\|^2.
\end{equation}
Summing with weights $p_k$ and invoking Assumption \ref{ass:data-heterogeneity} yields
\begin{align}
\label{eq:bound_drift_sum}
\sum_{k=1}^K p_k \mathbb{E}_t\|g_{t,i}^k\|^2
&\le \underbrace{2\sigma_l^2 + 4\kappa^2}_{C_\sigma} + 4M\|\nabla F(\omega_t)\|^2 \notag \\
&+ 4L^2 \sum_{k=1}^K p_k \mathbb{E}_t\|\omega_{t,i}^k - \omega_t\|^2.
\end{align}
Substituting (\ref{eq:bound_drift_sum}) into (\ref{eq:bound_drift}) and defining $D_e$ gives
\begin{equation}
D_e \le \eta_t^2 e \sum_{i=0}^{e-1} \left( C_\sigma + 4M\|\nabla F(\omega_t)\|^2 + 4L^2 D_i \right).
\end{equation}
For sufficiently small stepsize ($4L^2\eta_t^2 E^2 \le 1/2$), unrolling the recurrence and absorbing higher-order terms implies
\begin{equation}
D_e \le \mathcal{O}(\eta_t^2 E^2) \left( C_\sigma + \|\nabla F(\omega_t)\|^2 \right).
\end{equation}
For $\delta_t = \frac{1}{E} \sum_{e=0}^{E-1} \sum_{k=1}^K p_k (\nabla F_k(\omega_{t,e}^k) - \nabla F_k(\omega_t))$, Jensen's inequality and $L$-smoothness yield
\begin{equation}
\mathbb{E}_t\|\delta_t\|^2\!\le\!\frac{L^2}{E}\sum_{e=0}^{E-1}D_e\!\le\!\mathcal{O}\!\left(\!L^2\eta_t^2E^2\!\right)\!\left(\!C_\sigma\!+\!4M\|\nabla F(\omega_t)\|^2\!\right).
\end{equation}
\end{proof}

\begin{lemma}[Bounding the Ideal Gradient Variance]
Under Assumptions 1--3, the second moment of the ideal aggregated gradient $g_t$ is bounded by:
\begin{equation}
\mathbb{E}_t\|g_t\|^2 \le \big(1 + \mathcal{O}(L^2 \eta_t^2 E^2)\big) \cdot \Big( C_\sigma + 4M\|\nabla F(\omega_t)\|^2 \Big).
\end{equation}
\end{lemma}

\begin{proof}
By definition of $g_t$ and Jensen's inequality,
\begin{align}
\mathbb{E}_t\|g_t\|^2 &= \mathbb{E}_t\left\| \sum_{k=1}^K p_k \left( \frac{1}{E} \sum_{e=0}^{E-1} g_{t,e}^k \right) \right\|^2 \notag \\
&\le \frac{1}{E} \sum_{e=0}^{E-1} \sum_{k=1}^K p_k \mathbb{E}_t\|g_{t,e}^k\|^2.
\end{align}
Using variance decomposition,
\begin{equation}
\mathbb{E}_t\|g_{t,e}^k\|^2 \le \sigma_l^2 + \mathbb{E}_t\|\nabla F_k(\omega_{t,e}^k)\|^2.
\end{equation}
Adding and subtracting $\nabla F_k(\omega_t)$ and applying $L$-smoothness,
\begin{equation}
\mathbb{E}_t\|g_{t,e}^k\|^2 \le \sigma_l^2 + 2L^2\mathbb{E}_t\|\omega_{t,e}^k - \omega_t\|^2 + 2\|\nabla F_k(\omega_t)\|^2.
\end{equation}
Taking weighted sums and using Assumption 3,
\begin{equation}
\sum_{k=1}^K p_k \mathbb{E}_t\|g_{t,e}^k\|^2
\le \sigma_l^2 + 2\kappa^2 + 2M\|\nabla F(\omega_t)\|^2 + 2L^2 D_e.
\end{equation}
Averaging over $e$ and substituting Lemma 1,
\begin{align}
\mathbb{E}_t\|g_t\|^2 &\le \frac{1}{2}C_\sigma + 2M\|\nabla F(\omega_t)\|^2 + \frac{2L^2}{E} \sum_{e=0}^{E-1} D_e \\
&\le \frac{1}{2}C_\sigma + 2M\|\nabla F(\omega_t)\|^2 \notag \\
&+ \mathcal{O}(L^2 \eta_t^2 E^2) \left( C_\sigma + 4M\|\nabla F(\omega_t)\|^2 \right).
\end{align}
Regrouping terms gives the stated bound.
\end{proof}

\begin{lemma}[One-step Descent Lemma]
Under Assumptions 1--5, provided $\eta_t$ is sufficiently small, the expected objective value at iteration $t+1$ satisfies:
\begin{align}
\mathbb{E}_t[F(\omega_{t+1})] &\le F(\omega_t) - \mathcal{C}_0 \eta_t E \|\nabla F(\omega_t)\|^2 + \mathcal{C}_1 \eta_t^2 E^2 \notag \\
& + \mathcal{C}_2 \eta_t E B_\mu^2 + \mathcal{C}_3 \eta_t^2 E^2 B_V^2 + \tilde{R}_t,
\end{align}
where $\mathcal{C}_0, \mathcal{C}_1, \mathcal{C}_2, \mathcal{C}_3$ are positive constants, and $\tilde{R}_t = \left(\frac{8}{\eta_t E} + L\right)\|r_t\|^2$.
\end{lemma}

\begin{proof}
By Assumption 1 ($L$-smoothness),
\begin{equation}
F(\omega_{t+1}) \le F(\omega_t) + \langle \nabla F(\omega_t), \Delta_t \rangle + \frac{L}{2}\|\Delta_t\|^2.
\end{equation}
Substitute $\Delta_t = -\eta_t E(g_t + b_t) + r_t$ and decompose
\begin{align}
\mathbb{E}_t\langle \nabla F(\omega_t), \Delta_t \rangle
&= \underbrace{-\eta_t E \langle \nabla F, \mathbb{E}_t[g_t] \rangle}_{(A)}
+ \underbrace{-\eta_t E \mathbb{E}_t \langle \nabla F, b_t \rangle}_{(B)} \notag \\
&+ \underbrace{\mathbb{E}_t \langle \nabla F, r_t \rangle}_{(C)}.
\end{align}

Using $\mathbb{E}_t[g_t] = \nabla F + \delta_t$ and Young's inequality,
\begin{align}
(A) &= -\eta_t E \|\nabla F\|^2 - \eta_t E \langle \nabla F, \delta_t \rangle \notag \\
&\le -\frac{3\eta_t E}{4}\|\nabla F\|^2 + \eta_t E \|\delta_t\|^2.
\end{align}
Using $\mathbb{E}_{S_t}[b_t] = \mu_t$,
\begin{equation}
(B) \le \frac{\eta_t E}{8} \|\nabla F\|^2 + 2\eta_t E \|\mu_t\|^2 \le \frac{\eta_t E}{8} \|\nabla F\|^2 + 2\eta_t E B_\mu^2.
\end{equation}
And for $(C)$,
\begin{equation}
(C) \le \frac{\eta_t E}{16} \|\nabla F\|^2 + \frac{8}{\eta_t E}\|r_t\|^2.
\end{equation}

For the quadratic term, using $\|x+y+z\|^2 \le 3\|x\|^2 + 3\|y\|^2 + 3\|z\|^2$ and Assumption 4,
\begin{equation}
\frac{L}{2} \mathbb{E}_t\|\Delta_t\|^2 \le 2L\eta_t^2 E^2 \mathbb{E}_t\|g_t\|^2 + 2L\eta_t^2 E^2 (B_\mu^2 + B_V^2) + L\|r_t\|^2.
\end{equation}
Combining all bounds, the coefficient of $\|\nabla F(\omega_t)\|^2$ is
\begin{equation}
\text{Coeff} = -\eta_t E \left( 1 - \frac{1}{4} - \frac{1}{8} - \frac{1}{16} \right) + \text{Positive Erosion Terms.}
\end{equation}
By Lemmas 1--2, the positive erosion terms are of orders $\mathcal{O}(\eta_t^3)$ and $\mathcal{O}(\eta_t^4)$. For sufficiently small $\eta_t$ (e.g., $\eta_tEL\le c$), they are dominated by the first-order descent term, yielding a net coefficient $-\mathcal{C}_0\eta_tE$. Grouping the remaining constants gives the stated inequality.
\end{proof}

\subsection{Proof of Theorem 1}
\begin{proof}
Taking total expectation in Lemma 3 and summing over $t=0,\dots,T-1$,
\begin{align}
\mathcal{C}_0 E \sum_{t=0}^{T-1} \eta_t \mathbb{E}\|\nabla F(\omega_t)\|^2
&\le F(\omega_0) \!- \!\mathbb{E}[F(\omega_T)] \!+ \!\mathcal{C}_1 E^2 \sum_{t=0}^{T-1} \eta_t^2 \notag \\
& + \mathcal{C}_2 E B_\mu^2 \sum_{t=0}^{T-1} \eta_t + \mathcal{C}_3 E^2 B_V^2 \sum_{t=0}^{T-1} \eta_t^2 \notag \\
& + \sum_{t=0}^{T-1} \mathbb{E}[\tilde{R}_t].
\end{align}
Using $\mathbb{E}[F(\omega_T)] \ge F_{\inf}$ and $\Delta_F = F(\omega_0)-F_{\inf}$, we obtain
\begin{align}
\min_{0 \le t < T} \mathbb{E}\|\nabla F(\omega_t)\|^2
&\le \!\frac{1}{\mathcal{C}_0 E \sum \eta_t}
\Big[\sum \mathbb{E}[\tilde{R}_t]\! + \!\mathcal{C}_2 E B_\mu^2 \sum \eta_t  \notag \\
&\!+\! (\mathcal{C}_1 E^2 \!+\! \mathcal{C}_3 E^2 B_V^2)\sum \eta_t^2\!+\! \Delta_F\Big].
\end{align}
With $\eta_t = \frac{\eta_0}{\sqrt{t+1}}$,
\begin{enumerate}
\item \textbf{First-order sum:} $\sum_{t=0}^{T-1} \eta_t = \Theta(\sqrt{T})$.
\item \textbf{Second-order sum:} $\sum_{t=0}^{T-1} \eta_t^2 = \Theta(\ln T)$.
\item \textbf{Residual series:} Since $\tilde{R}_t = \left(\frac{8}{\eta_t E} + L\right)\|r_t\|^2$, and using Assumption 5 ($\|r_t\|^2 \le \mathcal{O}(t^{-2-2\rho})$):
\begin{equation}
\mathbb{E}[\tilde{R}_t] \propto \sqrt{t} \cdot t^{-2-2\rho} = t^{-1.5-2\rho}.
\end{equation}
Since $-1.5-2\rho<-1$, the $p$-series converges, i.e., $\sum_{t=0}^\infty \mathbb{E}[\tilde{R}_t] \le C_R < \infty$.
\end{enumerate}

Dividing each term on the right-hand side by $\sum \eta_t = \Theta(\sqrt{T})$:
\begin{itemize}
\item Initialization gap: $\frac{\Delta_F}{\Theta(\sqrt{T})} = \mathcal{O}\left(\frac{1}{\sqrt{T}}\right)$.
\item Variance and noise terms: $\frac{\Theta(\ln T)}{\Theta(\sqrt{T})} = \mathcal{O}\left(\frac{\ln T}{\sqrt{T}}\right)$.
\item Proxy residual: $\frac{C_R}{\Theta(\sqrt{T})} = \mathcal{O}\left(\frac{1}{\sqrt{T}}\right)$.
\item Systematic bias: $\frac{\mathcal{C}_2 E B_\mu^2 \Theta(\sqrt{T})}{\mathcal{C}_0 E \Theta(\sqrt{T})} = \mathcal{O}(B_\mu^2)$.
\end{itemize}

Combining these terms gives Theorem 1.
\end{proof}
\subsection{Proof of Corollary 1}
\begin{proof}
Taking $T\to\infty$ in Theorem 1 gives
\begin{equation}
\begin{aligned}
\lim_{T \to \infty} \Big[\, &\mathcal{O}\left( \frac{\Delta_F}{E\sqrt{T}} \right)
+ \mathcal{O}\left( \frac{(\sigma_l^2 + \kappa^2 + B_V^2) E \ln T}{\sqrt{T}} \right) \\
&+ \mathcal{O}\left( \frac{C_R}{E\sqrt{T}} \right)
+ \mathcal{O}\left( B_\mu^2 \right)\, \Big].
\end{aligned}
\end{equation}

Using standard limits,
$\lim_{T \to \infty} \frac{1}{\sqrt{T}} = 0$,
$\lim_{T \to \infty} \frac{\ln T}{\sqrt{T}} = 0$.
All terms that depend on $T$ vanish, and only the constant bias term remains:
\begin{equation}
\lim_{T \to \infty} \min_{0 \le t < T} \mathbb{E}\|\nabla F(\omega_t)\|^2 \le \mathcal{O}(B_\mu^2).
\end{equation}
thus establishing Corollary 1.
\end{proof}
%%%%%%%%%%%%%%%%%%%%%%%%%%%%%%%%%%%%%%%%%%%%%%%%%%%%%%%%%%%%%%

\bibliographystyle{IEEEtran}
\bibliography{bibfile}

%%%%%%%%%%%%%%%%%%%%%%%%%%%%%%%%%%%%%%%%%%%%%%%%%%%%%%%%%%%%%%

\begin{IEEEbiography}[{\includegraphics[width=1in,height=1.25in,clip,keepaspectratio]{fig/picture/Cuiyongsheng.jpg}}]{Yongsheng Cui}
	(Graduate Student Member, IEEE)
	is currently pursuing the Ph.D. degree in the Department of Electronic Engineering, Tsinghua University, Beijing, China.
	His research interests include deep space communications, satellite constellation design, 
	machine learning and satellite TT\&C systems.
\end{IEEEbiography}

\begin{IEEEbiography}[{\includegraphics[width=1in,height=1.25in,clip,keepaspectratio]{fig/picture/Chentaiyi.jpg}}]{Taiyi Chen}
	(Student Member, IEEE)
	received the Master degree in telecommunication engineering from the Department of Electronic Engineering, 
	Tsinghua University, Beijing, P.R. China. 
	His research interests include Wireless Communication Systems for satellite constellations, 
	channel coding, and deep space communications.
\end{IEEEbiography}

\begin{IEEEbiography}[{\includegraphics[width=1in,height=1.25in,clip,keepaspectratio]{fig/picture/Haoran Xie-eps-converted-to.pdf}}]{Haoran Xie}
	(Member, IEEE) 
    is a researcher at the Electric Power Research Institute, China Southern Power Grid.
    His research interests include application of optimization, machine learning, 
	and statistical theory. 
	He has served as the Chair for IEEE IWCMC 2022/2023/2024 Symposium and IEEE VTC 2022-Fall Symposium.
\end{IEEEbiography}

\begin{IEEEbiography}[{\includegraphics[width=1in,height=1.25in,clip,keepaspectratio]{fig/picture/YafengZhan-eps-converted-to.pdf}}]{Yafeng Zhan}
	(Member, IEEE)
	received B.S.E.E. and Ph.D.E.E. degrees from the Department of Electronic Engineering, Tsinghua University, Beijing, China, 
	in 1999 and 2004, respectively. 
	He is currently a Professor with the Beijing National Research Center for Information Science and Technology, 
	Tsinghua University. His current research interests include satellite TT\&C systems, 
	communication signal processing and deep space communications. 
    He serves as an editor of \textit{Chinese Space Science and Technology} and the \textit{Journal of Deep Space Exploration}.
\end{IEEEbiography}
\end{document}


